\documentclass[10pt,a4paper]{amsart}
\usepackage{amsmath,amssymb,amsthm,graphicx}
\usepackage{amsmath}
\usepackage{amsfonts}
\usepackage{amssymb}
\usepackage{graphicx}
\usepackage{float}
\usepackage{caption}
\usepackage{pgf,tikz}
\usetikzlibrary{arrows}
\usepackage{hyperref}

\definecolor{ffttww}{rgb}{1,0.2,0.4}
\definecolor{uuuuuu}{rgb}{0.27,0.27,0.27}
\definecolor{cqcqcq}{rgb}{0.75,0.75,0.75}
\definecolor{wwqqcc}{rgb}{0.4,0,0.8}

\newcommand{\R}{\mathbb{R}}
\newcommand{\N}{\mathbb{N}}
\newcommand{\HH}{\mathcal{H}}
\newcommand{\RR}{\mathcal{R}}
\newcommand{\A}{\mathcal{A}}
\newcommand{\B}{\mathcal{B}}
\newcommand{\est}{\emptyset}
\newtheorem{theorem}{Theorem}
\newtheorem{proposition}{Proposition}
\newtheorem{corollary}{Corollary}
\newtheorem{lemma}{Lemma}
\newtheorem{definition}{Definition}
\newtheorem{jel}{Notation}
\newtheorem{conj}{Conjecture}
\title{Uniform bottomless-rectangle hypergraphs with no3-shallow hitting set}
\author{Balázs Bursics, Bence Csonka, Luca Szepessy}
\date{arxiv:2307.12154v3; CC BY 4.0}
\begin{document}
\maketitle
\noindent\textit{Editorial note.} The counted hyperedges are captured m-point traces of bottomless rectangles, as defined at182-183 and used at470/474; source466 contained means captured in this range-family sense, not every subset of a larger trace.\par\medskip
\noindent\textit{Editorial note.} The word northwest for B\_i at464 conflicts with the same paragraph precise inequalities x3<x1<x2 and with the native figure. Those explicit coordinates place B\_i to the right of A\_i and C\_i to the left. The author text and figure are retained.\par\medskip
\noindent\textit{Editorial note.} The groups C\_i,j lie above b\_i\textasciicircum{}j and below b\_i\textasciicircum{}(j+1) when j<m; for j=m the precise condition y(b\_i\textasciicircum{}k)<y(C\_i,j) iff j>=k puts the final group above all B-points. No nonexistent b\_i\textasciicircum{}(m+1) is needed.\par\medskip
\noindent\textit{Editorial note.} The finite copies above the base points can be horizontally scaled into disjoint small neighborhoods, with every non-anchor point above the base line. This realizes the author separate copies without allowing a witnessing rectangle to capture points of another gadget.\par\medskip
\noindent\textit{Editorial note.} In Case2, the m-3 rightmost A-points exclude the leftmost anchor p\_i. The four disjoint consecutive triples in B\_i use the first12 B-points when m>12.\par\medskip
\noindent\textit{Editorial scope.} One lower-bound construction for shallow hitting sets in a uniform planar range hypergraph. The native figure is part of the proof. Arithmetic progression colorings, strip hypergraphs, and polychromatic upper bounds are not counted.\par\medskip
\noindent\textit{Allowed standard prerequisites (editorial).}\begin{itemize}
\item Finite planar range hypergraphs. A range captures exactly the input points lying in it. Bottomless rectangles have two vertical sides and a horizontal upper side, with no lower bound.. Verify each explicitly chosen edge has exactly m points.
\item Elementary counting and geometric separation. Use strict coordinate inequalities to set rectangle boundaries between finite point coordinates, and use four disjoint triples among at least12 ordered points.
\end{itemize}\medskip
\begin{definition} Denote by $H_{\ge m}$ or simply $H_m$ the hypergraph obtained from $H$ by deleting all hyperedges of size smaller than $m$, and denote by $H_{=m}$ the hypergraph consisting of the hyperedges of $H$ with size exactly $m$. Similarly, for a hypergraph family $\mathcal{H}$ let $\HH_m=\HH_{\ge m}=\{H_{\ge m}:H\in\HH\}$ and $\HH_{=m}=\{H_{=m}:H\in\HH\}$.
\end{definition}

\begin{definition}
 Let $H=(V,E)$ be a hypergraph, a $U\subset V$ vertex set $U\subset V$ is a {\it c-shallow hitting set} for a positive integer $c$ if for every hyperedge $e\in E$: $1\le|e\cap U|\le c.$
\end{definition} 

\begin{definition} Denote by $\mathcal{B}$ the hypergraph family which consists of those hypergraphs $H=(V,E)$, for which $V\subset \R^2$ is finite, and the edges are induced by bottomless rectangles: every $e\in E$ edge can be written in the form $e=\{(x,y)\in V: x_0 < x < x_1, y<y_0\}$ for some $x_0, x_1, y_0 \in\R$. 
\end{definition}

\noindent{\it Notation.}
Directions in the plane are denoted by north, west, south, and east.

\begin{theorem}\label{bottomless}
Let $m \ge 12$, then there is a member of $\mathcal{B}_{=m}$ which does not have a 3-shallow hitting set.
\end{theorem}

\begin{proof} [Proof of Theorem~\ref{bottomless}] 

Take $\B_{=m}$ with $m\ge 12$. We will show that this family has an element with no 3-shallow hitting set.
First take $m$ points, $P=\{p_1,\ldots,p_m\}$ along a horizontal line, and above each of them, take the construction on Figure \ref{fig:haromsek}.

%\begin{figure}
%    \centering
%    \includegraphics{pontoscucc12.png}
%    \caption{}
%    \label{fig:haromsek}
%\end{figure}

\begin{figure}[h!]
    \centering
    \begin{tikzpicture}[line cap=round,line join=round,>=triangle 45,x=1.1cm,y=1.1cm]
\clip(-4.02,-0.38) rectangle (7.54,7.76);
\draw [line width=0.6pt,dash pattern=on 2pt off 3pt,domain=-4.02:7.54] plot(\x,{(--13.9-0*\x)/3.5});
\draw [line width=0.6pt,dash pattern=on 2pt off 3pt,domain=-4.02:7.54] plot(\x,{(-42.49-0*\x)/-9.5});
\draw [line width=1pt,color=ffttww] (0,0) circle (0.24cm);
\draw (1.75,1.33) node[anchor=north west] [scale=0.8]{$A_i$};
\draw (-3.05,4.47) node[anchor=north west] [scale=0.8]{$C_{i,1}$};
\draw (-2.58,4.95) node[anchor=north west] [scale=0.8]{$C_{i,2}$};
\draw (-0.09,7.45) node[anchor=north west] [scale=0.8]{$C_{i,m}$};
\draw (5.68,5.98) node[anchor=north west] [scale=0.8]{$B_i$ };
\draw (-0.2,0) node[anchor=south east] 
[scale=0.8]{$P_i$ };
\draw [line width=0.6pt,dash pattern=on 2pt off 3pt,domain=-4.02:7.54] plot(\x,{(-47.24-0*\x)/-9.5});
\draw [line width=0.6pt,dash pattern=on 2pt off 3pt,domain=-4.02:7.54] plot(\x,{(-66.24-0*\x)/-9.5});
\begin{scriptsize}
\fill [color=black] (0,0) circle (2pt);
\fill [color=black] (1,1) circle (2pt);
\fill [color=black] (0.5,0.5) circle (2pt);
\fill [color=black] (6.5,4.47) circle (2pt);
\fill [color=black] (-3.47,4) circle (1.0pt);
\fill [color=black] (3,3) circle (2pt);
\fill [color=black] (7,3.97) circle (2pt);
\fill [color=black] (-3.25,4.22) circle (1.0pt);
\fill [color=black] (-3.03,4.44) circle (1.0pt);
\fill [color=black] (2,2) circle (2pt);
\fill [color=black] (1.5,1.5) circle (2pt);
\fill [color=black] (2.5,2.5) circle (1.0pt);
\fill [color=black] (2.75,2.75) circle (1.0pt);
\fill [color=black] (2.25,2.25) circle (1.0pt);
\fill [color=black] (4,6.97) circle (2pt);
\fill [color=black] (4.25,6.72) circle (1.0pt);
\fill [color=black] (4.5,6.47) circle (1.0pt);
\fill [color=black] (4.75,6.22) circle (1.0pt);
\fill [color=black] (5,5.97) circle (2pt);
\fill [color=black] (5.5,5.47) circle (2pt);
\fill [color=black] (6,4.97) circle (2pt);
\fill [color=black] (-3.4,4.07) circle (1.0pt);
\fill [color=black] (-3.32,4.15) circle (1.0pt);
\fill [color=black] (-3.18,4.3) circle (1.0pt);
\fill [color=uuuuuu] (-3.1,4.37) circle (0.5pt);
\fill [color=uuuuuu] (-3.14,4.33) circle (0.5pt);
\fill [color=black] (-3.1,4.37) circle (0.5pt);
\fill [color=black] (-3.06,4.41) circle (0.5pt);
\fill [color=black] (-3.14,4.33) circle (0.5pt);
\fill [color=black] (-2.97,4.5) circle (1.0pt);
\fill [color=black] (-2.9,4.57) circle (1.0pt);
\fill [color=black] (-2.82,4.65) circle (1.0pt);
\fill [color=black] (-2.75,4.72) circle (1.0pt);
\fill [color=black] (-2.68,4.8) circle (1.0pt);
\fill [color=uuuuuu] (-2.53,4.94) circle (1.0pt);
\fill [color=uuuuuu] (-2.6,4.87) circle (0.5pt);
\fill [color=uuuuuu] (-2.64,4.83) circle (0.5pt);
\fill [color=black] (-2.6,4.87) circle (0.5pt);
\fill [color=black] (-2.56,4.91) circle (0.5pt);
\fill [color=uuuuuu] (-2.43,5.04) circle (0.5pt);
\fill [color=black] (-2.34,5.13) circle (0.5pt);
\fill [color=uuuuuu] (-2.25,5.22) circle (0.5pt);
\fill [color=black] (-3.32,4.15) circle (1.0pt);
\fill [color=black] (-0.47,7) circle (1.0pt);
\fill [color=black] (-0.4,7.07) circle (1.0pt);
\fill [color=black] (-0.32,7.15) circle (1.0pt);
\fill [color=black] (-0.25,7.22) circle (1.0pt);
\fill [color=uuuuuu] (-0.14,7.33) circle (0.5pt);
\fill [color=black] (-0.18,7.3) circle (1.0pt);
\fill [color=uuuuuu] (-0.06,7.41) circle (0.5pt);
\fill [color=black] (-0.1,7.37) circle (0.5pt);
\fill [color=uuuuuu] (-0.03,7.44) circle (1.0pt);
\end{scriptsize}
\end{tikzpicture}
    \caption{Construction for Theorem~\ref{bottomless}.}
    \label{fig:haromsek}
\end{figure}


Starting from $p_i$ we have $A_i$ on a diagonal line consisting of $m-2$ points (including $p_i$) with increasing $x$ and $y$ coordinates, to northwest from $A_i$ we have $B_i=\{b_i^1,\ldots,b_i^m\}$ on a diagonal line with increasing $y$, but decreasing $x$ coordinates. To the west from $A_i$, we have $m^2$ points in a diagonal again with increasing $x$ and $y$ coordinates, consisting of $m$ smaller groups $C_{i,j}=\{c_{i,j}^1,\ldots,c_{i,j}^m\}$ (for $j=1,\ldots,m$) of size $m$, with $C_{i,j}$ being vertically between the $b_i^j$ and $b_i^{j+1}$. More precisely, let $(x_1,y_1)\in A_i,$ $(x_2,y_2)=b_i^k\in B_i,$ and $(x_3,y_3)\in C_{i,j}.$ Then $x_3<x_1<x_2$, $y_1<y_2$, $y_1<y_3,$ and we have $y_2<y_3$ exactly when $j\ge k.$

Let $H$ be the hypergraph on this set which contains all $m$-sets contained in a bottomless rectangle as hyperedges. Now for the sake of contradiction suppose that this hypergraph has a 3-shallow hitting set.

There is a point in $P$ which is in the hitting set, because these vertices form a hyperedge. Now we separate two cases:

Case 1. $A_i$ contains at least one further point from the hitting set besides $P$. $B_i$ consists of $m$ points, and can be separated by a bottomless rectangle, so it has a point from the hitting set, say $b_i^j$. We can choose an axis-parallel rectangle that contains exactly $b_i^j$ from $B_i$. Also in $C_{i,j}$ there must be a point from the hitting set, so we can choose the rectangle in a way that its top left corner contains exactly this point from $C_{i,j}$, say $c_{i,j}^l$. Now we got a bottomless rectangle with $m$ points in it, $m-2$ from $A_i$ with at least two points from the hitting set, and $b_i^j$ and $c_{i,j}^l$, both latter two from the hitting set. This is 4 in total; thus, the hitting set is not 3-shallow, a contradiction.



Case 2. $A_i$ does not contain any more points from the hitting set (apart from $p_i$). Now we can take a bottomless rectangle which contains the $m-3$ rightmost points of $A_i$ and any 3 neighbouring points from $B_i$. This implies that any 3 neighbouring points in $B_i$ contain a point from the hitting set, which means $B_i$ contains at least 4 points from the hitting set if $m\ge12$, which is a contradiction.
\end{proof}

\end{document}
